\subsection{Comparación global de estrategias}

% La lectura conjunta de los resultados de Grover y QFT permite extraer dos conclusiones
% generales. La primera es que la utilidad de la compresión depende de la estructura del estado
% cuántico y no puede asumirse de manera uniforme para cualquier carga de trabajo. La segunda
% es que, dentro de los escenarios favorables observados en esta campaña, Blosc ofrece la
% relación más equilibrada entre ahorro de memoria, tiempo de ejecución y exactitud numérica,
% mientras que ZFP solo resulta competitivo en casos específicos y a costa de introducir una
% discrepancia no nula respecto a la referencia.

\begin{table}[H]
    \centering
    \small
    \caption{Síntesis comparativa frente al almacenamiento no comprimido}
    \label{tab:summary-results}
    \begin{tabular}{>{\raggedright\arraybackslash}p{3.3cm}>{\centering\arraybackslash}p{1.8cm}>{\centering\arraybackslash}p{1.8cm}>{\centering\arraybackslash}p{1.8cm}>{\centering\arraybackslash}p{1.8cm}>{\centering\arraybackslash}p{2cm}}
        \hline
        \textbf{Carga de trabajo} & \textbf{Blosc: ahorro (\%)} & \textbf{Blosc: tiempo (x)} & \textbf{ZFP: ahorro (\%)} & \textbf{ZFP: tiempo (x)} & \textbf{ZFP: error} \\
        \hline
        Grover, objetivo único & 77.8 a 93.5 & 4.22 a 4.85 & 36.4 a 45.7 & 32.27 a 37.78 & $8.04\times10^{-15}$ a $1.61\times10^{-14}$ \\
        Grover, multiobjetivo & 77.8 a 93.4 & 4.17 a 4.86 & 36.7 a 45.7 & 33.12 a 38.13 & $8.04\times10^{-15}$ a $1.61\times10^{-14}$ \\
        QFT, superposición periódica & 51.4 a 75.3 & 2.54 a 3.48 & 74.2 a 83.2 & 4.41 a 5.75 & $3.79\times10^{-11}$ \\
        QFT, alta entropía & -49.0 a -12.1 & 10.24 a 12.52 & -55.5 a -2.1 & 85.65 a 103.01 & $2.82\times10^{-11}$ a $3.79\times10^{-11}$ \\
        \hline
    \end{tabular}
\end{table}

% Grover constituye el escenario más claramente favorable para la compresión sin pérdida. En
% ambas variantes evaluadas, Blosc supera el 77\% de ahorro de memoria en todos los tamaños y
% alcanza más de 93\% en 24 qubits, con un costo temporal cercano a $4\times$. ZFP, aunque
% mantiene errores máximos absolutos muy pequeños, del orden de $10^{-14}$, reduce menos
% memoria y exige sobrecostos temporales superiores a $30\times$. En este contexto, la
% estrategia sin pérdida no solo preserva la exactitud exacta, sino que también ofrece el mejor
% balance práctico.

% QFT obliga a matizar esa lectura. Cuando la entrada es periódica, tanto Blosc como ZFP
% resultan útiles, pero con perfiles distintos: Blosc ofrece una combinación más conservadora
% de tiempo, memoria y exactitud, mientras que ZFP logra el mayor ahorro espacial a costa de
% una mayor penalización temporal y de errores del orden de $10^{-11}$. Sin embargo, cuando la
% entrada presenta alta entropía, ambas estrategias dejan de ser ventajosas: el ahorro de
% memoria se vuelve negativo y el tiempo de ejecución aumenta de forma marcada, incluso en ZFP,
% a pesar de que el error admitido sigue siendo pequeño.

En conjunto, la evidencia experimental delimita con claridad el alcance de la propuesta. La
compresión sin pérdida del vector de estado es una alternativa viable cuando la memoria es la
restricción dominante y el estado conserva regularidades aprovechables. No obstante, sus
beneficios no son universales: dependen de la estructura concreta del estado y del patrón de
acceso que induce el circuito. Bajo las condiciones evaluadas, Blosc emerge como la opción
más robusta entre las estrategias comparadas, mientras que la compresión con pérdida basada
en ZFP solo ofrece ventajas parciales y situacionales.
